\cotefile{90}
\scanpage{14}
\deleted{\uline{Proposition} Soit $A$ anneau de Dedekind, $G$ gpe
\inserted{plat} affine \inserted{de t.f.} sur $A$, d'algèbre $B$ satisfaisant aux conditions
du Lemme 1. \deleted{Alors} Soit $N$ comme dans le Lemme 2,
et \deleted{\unsure{la}} considérons la représentation de $G$ dans $N$ induite
par la \unsure{repr.|représ.} régulière gauche. Alors l'hom.
$G\longrightarrow\underline{\mathrm{Aut}}(N)$ est un morphisme \unsure{fidèlem$^{\mathrm t}$|fidèlement} plat}

\[
\left\{\begin{aligned}
 &\underline{B}=\varinjlim B_i\\
1)\quad &B/B_i\ \text{plat}\quad\text{\deleted{\unsure{donc|dans}}}\qquad B_i\otimes B_i\subset B\otimes B\\
2)\quad &\Delta B_i\subset B_i\otimes B_i\\
3)\quad &B_i\ \text{loc libre de rt. f.}\quad \deleted{\text{\unsure{Suppos.|Supposons} $X$ qu. cpt qu. sép.}}
\end{aligned}\right.
\]

Soit $E$ un $\underline{O}_{X}G$-module qu. coh\unsure{.}\ \deleted{\illegible{2}}

Alors 1) Pour tt sous-module qu. coh. $F$
il existe un plus petit ss-module qu.
coh. $F'$ de $F$ \deleted{tel que} stable par $G$, et
sa formation commute à tt \unsure{changem$^{\mathrm t}$|changement} de
base plat. Si $F$ est de t.f. $F'$ itou

2) Si $X$ quasi-compact quasi-séparé, alors
$E$ est $\varinjlim$ de sous-modules \inserted{qu. coh.} de t.f.
qui sont stables par $G$.

3) \inserted{\deleted{\unsure{Si} $X$ \unsure{qu. cpt qu. sép.}} Si $X$ qu. cpt qu. sép.} Soit $F\unsure{\longrightarrow|-}E$ un hom, avec $F$ quasi-coh \inserted{de t.f.}
\unsure{ou|de} d'image \unsure{engendrant|engdr|engendre} $E$. Alors il existe
% archive sheet number: 13
